%% Chapitre2.tex — État de l'art
\chapter{État de l'art}
\label{chapter:etatart}

\epigraphe{Si j'ai vu plus loin, c'est en montant sur les épaules de géants.}{Isaac Newton}

%% ── Introduction ───────────────────────────────────────────────────────────
\section*{Introduction}
\addcontentsline{toc}{section}{Introduction}
\par Ce chapitre présente les fondements théoriques mobilisés par le projet. Il
ne s'agit pas seulement de définir quelques outils, mais de situer la solution
développée dans un cadre méthodologique cohérent : une chaîne décisionnelle
complète, un entrepôt de données structuré, des indicateurs financiers
normalisés et une démarche de prévision appuyée sur l'apprentissage automatique.

\par Le projet de la \sfbt combine en effet plusieurs domaines complémentaires.
L'informatique décisionnelle fournit l'architecture générale permettant de
transformer des fichiers hétérogènes en information exploitable. La modélisation
dimensionnelle organise les données dans une structure adaptée à l'analyse. La
finance d'entreprise fournit les ratios, les règles de comparaison et les
principes d'interprétation. Enfin, le \textit{Machine Learning} apporte une
dimension prédictive, utile pour dépasser le diagnostic historique et produire
des estimations futures.

\par Le chapitre est structuré comme suit. La première section présente les
principes de la \ac{BI} et de la chaîne \ac{ETL}. La deuxième traite de la
qualité des données, enjeu essentiel dans un contexte financier. La troisième
section expose les modèles d'entrepôts de données et justifie le choix du schéma
en étoile. La quatrième section développe l'analyse financière par les ratios et
le référentiel MSI20000. La cinquième présente le benchmark financier. La
dernière section introduit les modèles de prévision et les métriques
d'évaluation.

%% ───────────────────────────────────────────────────────────────────────────
\section{Informatique décisionnelle (Business Intelligence)}
\label{sec:bi}
\par L'informatique décisionnelle, ou \textit{Business Intelligence} (\ac{BI}),
désigne l'ensemble des méthodes, architectures et outils qui permettent de
collecter, consolider, analyser et restituer les données d'une organisation afin
d'améliorer la prise de décision. Elle transforme des données opérationnelles,
souvent dispersées et difficiles à lire, en indicateurs synthétiques, tableaux de
bord, rapports et analyses comparatives.

\par Dans une organisation, les données brutes ne deviennent pas immédiatement
une information utile. Elles doivent d'abord être extraites depuis leurs sources,
nettoyées, harmonisées, historisées puis restituées sous une forme lisible. Cette
logique est particulièrement importante dans le domaine financier, où une erreur
de libellé, une colonne comparative mal interprétée ou une devise non précisée
peut conduire à une conclusion incorrecte sur la santé d'une entreprise.

\subsection{Rôle de la BI dans le pilotage financier}
\par Le pilotage financier repose sur une lecture régulière des états
financiers, sur le suivi des écarts et sur la comparaison des résultats avec des
objectifs ou des références sectorielles. La \ac{BI} facilite ce pilotage en
automatisant la préparation des données et en donnant aux décideurs une vision
rapide de la situation. Elle permet notamment :
\begin{itemize}[noitemsep]
  \item de centraliser des données issues de sources multiples ;
  \item de fiabiliser les indicateurs utilisés dans les tableaux de bord ;
  \item de suivre l'évolution d'une entreprise sur plusieurs exercices ;
  \item de comparer plusieurs sociétés sur des ratios homogènes ;
  \item d'appuyer la décision par des visualisations claires et actualisables.
\end{itemize}

\par Dans le cas de la \sfbt, la BI joue un rôle structurant : elle permet de
passer de fichiers Excel bruts, hétérogènes et parfois incomplets, à un système
reproductible capable de calculer des ratios, d'alimenter un entrepôt de données
et de produire des tableaux de bord sous \textit{Power BI}.

\subsection{Architecture générale d'un système décisionnel}
\par Une architecture décisionnelle classique se compose de quatre couches :
les sources de données, la chaîne d'intégration, le stockage analytique et la
restitution. La figure~\ref{fig:bichain} représente cette organisation en
couches.

\begin{figure}[H]
\centering
\begin{tikzpicture}[
  node distance=0.6cm,
  couche/.style={rectangle, rounded corners, draw=darkblue, very thick,
    fill=lightblue!45, minimum width=11.5cm, minimum height=1.1cm,
    align=center, font=\small},
  arr/.style={-{Latex[length=2.5mm]}, very thick, darkblue}
]
\node[couche] (src) {\textbf{Sources de données}\\Fichiers Excel, CSV, états financiers, systèmes comptables};
\node[couche, below=of src] (etl) {\textbf{Intégration ETL}\\Extraction, nettoyage, normalisation, contrôle et historisation};
\node[couche, below=of etl] (dw) {\textbf{Stockage analytique}\\Data Warehouse, tables de faits, dimensions, vues métiers};
\node[couche, below=of dw] (visu) {\textbf{Restitution et analyse}\\Tableaux de bord, ratios, score, benchmark, prévision};
\draw[arr] (src) -- (etl);
\draw[arr] (etl) -- (dw);
\draw[arr] (dw) -- (visu);
\end{tikzpicture}
\caption[Architecture en couches d'un système décisionnel]{Architecture en couches d'un système décisionnel (BI).}
\label{fig:bichain}
\end{figure}

\par Cette architecture répond à une logique de séparation des responsabilités.
Les sources conservent les données d'origine ; la couche ETL assure la qualité et
l'harmonisation ; l'entrepôt organise les données selon les axes d'analyse ; les
outils de restitution rendent l'information accessible à l'utilisateur final. Ce
découpage rend la solution plus robuste, car une modification dans les sources
ou dans les tableaux de bord ne remet pas en cause toute l'architecture.

\subsection{Le processus ETL}
\par Le processus \ac{ETL} constitue le coeur technique d'un système
décisionnel. Il regroupe trois étapes principales : l'extraction
(\textit{Extract}), la transformation (\textit{Transform}) et le chargement
(\textit{Load}). Dans les projets modernes, cette logique peut être complétée par
des contrôles de qualité, des rapports d'anomalies et des scripts de
reproductibilité. Le tableau~\ref{tab:etl} détaille ces étapes et leur
application au projet.

\begin{table}[H]
\centering
\caption[Étapes du processus ETL]{Étapes du processus ETL et application au projet.}
\label{tab:etl}
\begin{tabularx}{\linewidth}{|p{3cm}|X|X|}
\hline
\textbf{Étape} & \textbf{Principe général} & \textbf{Application au projet \sfbt} \\ \hline
Extraction & Lire les données depuis les systèmes sources. & Lecture des fichiers Excel de la \sfbt, DELICE, AB~InBev et Coca-Cola. \\ \hline
Transformation & Nettoyer, normaliser, convertir et enrichir les données. & Normalisation des libellés, classification des états, correction des doublons, ajout de la devise et de l'échelle. \\ \hline
Chargement & Stocker les données prêtes pour l'analyse. & Génération des fichiers CSV/Parquet et alimentation du Data Warehouse SQLite. \\ \hline
Contrôle & Vérifier la cohérence des données produites. & Rapport qualité, contrôle Actif = Passif, collisions, tests automatiques. \\ \hline
\end{tabularx}
\end{table}

\par L'ETL ne doit donc pas être vu comme une simple conversion de format. Dans
un projet financier, il assure une fonction de fiabilisation métier : il protège
l'analyse contre les erreurs de saisie, les variations de libellés, les
différences de périodicité et les problèmes d'échelle.

%% ───────────────────────────────────────────────────────────────────────────
\section{Qualité et préparation des données financières}
\label{sec:dataquality}
\par La qualité des données est une condition préalable à toute analyse
financière fiable. Une donnée financière est sensible parce qu'elle supporte des
décisions importantes : évaluer la solvabilité, mesurer la rentabilité, comparer
une entreprise à ses concurrents ou anticiper ses résultats. Une valeur mal
classée peut modifier un ratio, un score ou une conclusion stratégique.

\subsection{Dimensions de la qualité des données}
\par La littérature distingue plusieurs dimensions de qualité : exactitude,
complétude, cohérence, unicité, fraîcheur et traçabilité. Dans un projet
d'analyse financière, ces dimensions prennent une importance particulière. Le
tableau~\ref{tab:qualite} les résume.

\begin{table}[H]
\centering
\caption[Dimensions de qualité des données]{Principales dimensions de qualité des données financières.}
\label{tab:qualite}
\begin{tabularx}{\linewidth}{|p{3.2cm}|X|X|}
\hline
\textbf{Dimension} & \textbf{Définition} & \textbf{Exemple dans ce projet} \\ \hline
Exactitude & Les valeurs correspondent à la réalité comptable. & Correction des colonnes comparatives recopiées par erreur. \\ \hline
Complétude & Les données nécessaires sont présentes. & Signalement des exercices annuels absents pour certaines années. \\ \hline
Cohérence & Les relations comptables sont respectées. & Contrôle de l'équilibre Actif = Passif. \\ \hline
Unicité & Une même information n'est pas comptée plusieurs fois. & Suppression des doublons strictement identiques. \\ \hline
Traçabilité & L'origine d'une valeur reste identifiable. & Conservation du libellé brut et du fichier source. \\ \hline
Comparabilité & Les indicateurs peuvent être comparés dans le temps et entre sociétés. & Ajout de la devise, de l'échelle et calcul de ratios sans dimension. \\ \hline
\end{tabularx}
\end{table}

\subsection{Normalisation des libellés}
\par Les états financiers contiennent souvent des libellés écrits de manière
différente selon les années : accents variables, majuscules, espaces
insécables, retours à la ligne, astérisques de notes, abréviations ou
reformulations. Pour un humain, ces variantes restent compréhensibles ; pour un
programme, elles peuvent être interprétées comme des indicateurs différents.

\par La normalisation des libellés vise donc à construire une clé stable pour
chaque concept financier. Cette clé sert à regrouper les séries temporelles et à
retrouver automatiquement les postes nécessaires au calcul des ratios. Dans le
cadre de ce projet, deux niveaux sont distingués : un libellé propre pour
l'affichage et une clé canonique pour les jointures. La
figure~\ref{fig:normalisation} illustre ce passage du libellé brut à la clé
stable.

\begin{figure}[H]
\centering
\begin{tikzpicture}[
  node distance=0.8cm and 1cm,
  box/.style={rectangle, rounded corners, draw=midblue, thick, fill=lightblue!35,
    align=center, minimum width=3cm, minimum height=1cm, font=\footnotesize},
  arr/.style={-{Latex[length=2.2mm]}, thick, midblue}
]
\node[box] (raw) {Libellé brut\\\textit{RESULTAT NET}\\\textit{DE LA PERIODE}};
\node[box, right=of raw] (clean) {Libellé propre\\Résultat net de la période};
\node[box, right=of clean] (key) {Clé canonique\\resultatnetdelexercice};
\draw[arr] (raw) -- (clean);
\draw[arr] (clean) -- (key);
\end{tikzpicture}
\caption[Normalisation des libellés financiers]{Principe de normalisation des libellés financiers.}
\label{fig:normalisation}
\end{figure}

\par Cette démarche est essentielle pour reconstituer des séries temporelles
longues. Sans normalisation, un même poste tel que le résultat net pourrait être
fragmenté en plusieurs indicateurs, empêchant le calcul fiable des tendances et
des ratios.

\subsection{Traitement des valeurs manquantes et des collisions}
\par Les valeurs manquantes et les collisions sont deux difficultés fréquentes
dans les données financières. Une valeur manquante peut provenir d'un poste non
publié, d'une cellule non numérique ou d'une année absente. Une collision
apparaît lorsqu'un même indicateur, pour une même société et une même date,
porte plusieurs valeurs distinctes. Elle peut traduire une ambiguïté métier, un
doublon source ou une absence de distinction entre états financiers.

\par Dans une démarche rigoureuse, ces cas ne doivent pas être masqués sans
justification. Le choix retenu est de supprimer uniquement les doublons exacts,
de corriger les erreurs identifiables et de conserver les collisions résiduelles
avec un indicateur d'occurrence. Cette stratégie maintient la traçabilité et
évite de prendre une décision automatique lorsqu'une validation métier est
nécessaire.

%% ───────────────────────────────────────────────────────────────────────────
\section{Entrepôts de données et modélisation dimensionnelle}
\label{sec:dw}
\par Un \textbf{entrepôt de données} est une base orientée sujet, intégrée,
historisée et non volatile, conçue pour l'analyse plutôt que pour la
production~\cite{inmon2005}. Contrairement aux bases opérationnelles, qui
optimisent la saisie et les transactions, l'entrepôt optimise les requêtes
analytiques, les agrégations et les croisements selon plusieurs axes.

\subsection{Approches Inmon et Kimball}
\par Deux approches classiques structurent la conception des entrepôts de
données. L'approche d'Inmon privilégie un entrepôt d'entreprise centralisé,
normalisé, construit comme une source unique de vérité. L'approche de Kimball
met l'accent sur les \textit{data marts} dimensionnels orientés métier,
connectés par des dimensions conformes~\cite{kimball2013}. Les deux approches ne
s'opposent pas totalement ; elles répondent à des contextes différents. Le
tableau~\ref{tab:inmonkimball} en propose une comparaison synthétique.

\begin{table}[H]
\centering
\caption[Comparaison Inmon et Kimball]{Comparaison synthétique des approches Inmon et Kimball.}
\label{tab:inmonkimball}
\begin{tabularx}{\linewidth}{|p{3cm}|X|X|}
\hline
\textbf{Critère} & \textbf{Approche Inmon} & \textbf{Approche Kimball} \\ \hline
Logique & Entrepôt central d'entreprise. & Data marts dimensionnels orientés métier. \\ \hline
Modélisation & Souvent normalisée. & Schéma en étoile ou constellation. \\ \hline
Déploiement & Plus long, fortement structurant. & Progressif, plus rapide pour les usages analytiques. \\ \hline
Utilisateur cible & Gouvernance globale de la donnée. & Analyse métier et tableaux de bord. \\ \hline
Pertinence pour ce projet & Utile comme principe d'intégration. & Très adaptée aux ratios, au temps, aux sociétés et à Power BI. \\ \hline
\end{tabularx}
\end{table}

\par Pour le présent projet, l'approche dimensionnelle est la plus pertinente :
les analyses se font principalement par société, période, indicateur et ratio.
Ces axes correspondent naturellement à des dimensions, tandis que les valeurs
financières et les ratios constituent des faits.

\subsection{Tables de faits et dimensions}
\par La modélisation dimensionnelle distingue les tables de faits et les tables
de dimensions. Une table de faits contient les mesures numériques observées :
montant d'un poste financier, valeur d'un ratio, note ou score. Les dimensions
décrivent le contexte de ces mesures : date, société, indicateur, ratio, secteur
ou devise.

\par Cette séparation facilite les analyses multidimensionnelles. Par exemple,
une même valeur de ratio peut être analysée selon l'année, la société, la
catégorie du ratio ou le type de période. La granularité doit être définie avec
soin : elle précise ce que représente exactement une ligne de la table de faits.
Dans ce projet, la granularité est fine : une valeur pour une société, une date
et un indicateur ou ratio.

\subsection{Schéma en étoile}
\par Le \textbf{schéma en étoile} organise une table de faits centrale entourée
de tables de dimensions. Sa simplicité le rend particulièrement adapté aux
outils de \textit{reporting} comme \textit{Power BI}, qui exploitent nativement
ce type de modèle. On le distingue du \textit{schéma en flocon}, dans lequel les
dimensions sont normalisées, et du \textit{schéma en constellation}, où plusieurs
tables de faits partagent des dimensions communes. La figure~\ref{fig:schemas}
représente le principe général du schéma en étoile.

\begin{figure}[H]
\centering
\begin{tikzpicture}[
  dim/.style={rectangle, rounded corners, draw=midblue, thick, fill=lightblue,
    align=center, minimum height=1cm, minimum width=2.8cm, font=\footnotesize},
  fact/.style={rectangle, draw=espritred, very thick, fill=espritred!10,
    align=center, minimum height=1.3cm, minimum width=3.4cm, font=\footnotesize\bfseries},
  link/.style={thick, midblue}
]
\node[fact] (f) {Table de faits\\Mesures financières};
\node[dim, above=1.3cm of f] (t) {Dim Temps};
\node[dim, below=1.3cm of f] (i) {Dim Indicateur};
\node[dim, left=2cm of f] (s) {Dim Société};
\node[dim, right=2cm of f] (r) {Dim Ratio};
\draw[link] (f) -- (t);
\draw[link] (f) -- (i);
\draw[link] (f) -- (s);
\draw[link] (f) -- (r);
\end{tikzpicture}
\caption[Principe du schéma en étoile]{Principe du schéma en étoile appliqué aux données financières.}
\label{fig:schemas}
\end{figure}

\par Le tableau~\ref{tab:modelesdimensionnels} compare les principaux modèles
dimensionnels et justifie le choix retenu.

\begin{table}[H]
\centering
\caption[Comparaison des modèles dimensionnels]{Comparaison des modèles dimensionnels classiques.}
\label{tab:modelesdimensionnels}
\begin{tabularx}{\linewidth}{|p{3.2cm}|X|X|}
\hline
\textbf{Modèle} & \textbf{Avantages} & \textbf{Limites} \\ \hline
Schéma en étoile & Simple, lisible, performant pour les tableaux de bord, facile à comprendre par les utilisateurs métier. & Peut introduire une certaine redondance dans les dimensions. \\ \hline
Schéma en flocon & Réduit la redondance et normalise davantage les dimensions. & Plus complexe à interroger et moins intuitif dans les outils BI. \\ \hline
Constellation & Permet plusieurs processus analytiques partageant des dimensions communes. & Demande une conception plus avancée et une gouvernance stricte des dimensions. \\ \hline
\end{tabularx}
\end{table}

\par Le choix du schéma en étoile se justifie donc par la nature du besoin : il
faut alimenter des dashboards, comparer des sociétés et suivre des ratios dans
le temps. La priorité est la lisibilité analytique plutôt que la normalisation
maximale.

%% ───────────────────────────────────────────────────────────────────────────
\section{Analyse financière par les ratios}
\label{sec:ratios}
\par L'analyse financière par les ratios~\cite{vernimmen2022} consiste à
rapprocher des postes des états financiers pour obtenir des indicateurs plus
interprétables que les montants bruts. Un montant isolé a souvent peu de sens :
un résultat net de plusieurs millions peut être excellent pour une petite
entreprise ou faible pour un groupe international. Un ratio, en revanche,
exprime une relation : résultat net sur chiffre d'affaires, dettes sur actifs,
actifs courants sur passifs courants, etc.

\subsection{Intérêt des ratios}
\par Les ratios présentent trois avantages majeurs. Premièrement, ils réduisent
l'effet de taille entre entreprises. Deuxièmement, ils facilitent le suivi dans
le temps. Troisièmement, ils permettent de comparer une entreprise à des normes,
à des concurrents ou à des benchmarks sectoriels. Cette logique est au coeur du
diagnostic financier.

\par Les ratios doivent toutefois être interprétés avec prudence. Ils dépendent
du secteur d'activité, des méthodes comptables, de la structure de l'entreprise
et de la qualité des données sources. Une entreprise industrielle, par exemple,
n'a pas les mêmes besoins d'immobilisations qu'une société de services ; une
holding peut présenter des revenus comptables très différents d'une société
opérationnelle.

\subsection{Familles de ratios}
\par Les ratios financiers se regroupent en grandes familles. Chaque famille
répond à une question de gestion : l'entreprise peut-elle payer ses dettes à
court terme ? Est-elle rentable ? Est-elle trop endettée ? Utilise-t-elle
efficacement ses actifs ? Le tableau~\ref{tab:famillesratios} synthétise les
principales familles retenues dans ce projet.

\begin{table}[H]
\centering
\caption[Principales familles de ratios]{Principales familles de ratios financiers.}
\label{tab:famillesratios}
\begin{tabularx}{\linewidth}{|p{3.2cm}|X|X|}
\hline
\textbf{Famille} & \textbf{Question analysée} & \textbf{Exemples d'indicateurs} \\ \hline
Liquidité & L'entreprise peut-elle couvrir ses engagements à court terme ? & Liquidité générale, liquidité réduite, liquidité immédiate. \\ \hline
Structure financière & Quelle est la part des capitaux propres et des dettes ? & Autonomie financière, endettement global, solvabilité. \\ \hline
Gestion des actifs & Les actifs sont-ils mobilisés efficacement ? & Levier économique, poids des immobilisations, vétusté. \\ \hline
Rentabilité & L'activité génère-t-elle un résultat satisfaisant ? & Marge nette, ROA, ROE, ROCE. \\ \hline
Cycle d'exploitation & Combien de temps l'argent reste-t-il immobilisé dans le cycle client-stock-fournisseur ? & DSO, DPO, DIO, BFR en jours. \\ \hline
Risque & L'entreprise présente-t-elle un risque financier ou commercial ? & Conan \& Holder, risque commercial, risque de placement. \\ \hline
\end{tabularx}
\end{table}

\subsection{Interprétation financière}
\par L'interprétation d'un ratio ne consiste pas uniquement à dire s'il est
haut ou bas. Elle exige de comprendre sa signification économique. Une liquidité
générale élevée peut signaler une bonne capacité de paiement, mais aussi un
excès de ressources non utilisées. Un endettement faible réduit le risque, mais
peut limiter l'effet de levier. Une marge nette élevée peut traduire une forte
rentabilité, mais doit être rapprochée du modèle économique de la société.

\par La lecture doit donc être croisée : liquidité, rentabilité, structure
financière et cycle d'exploitation doivent être analysés ensemble. C'est pour
cette raison qu'un score global peut être utile : il synthétise plusieurs
dimensions tout en conservant la possibilité de revenir au détail.

\subsection{Limites des ratios}
\par Malgré leur utilité, les ratios présentent des limites. Ils sont calculés à
partir d'états financiers historiques et ne captent pas toujours les événements
récents. Ils peuvent être sensibles aux choix comptables, aux éléments
exceptionnels et à la saisonnalité. Ils ne remplacent pas l'analyse qualitative :
position concurrentielle, stratégie, gouvernance, risques réglementaires et
contexte macroéconomique.

\par Le tableau~\ref{tab:limitesratios} résume les principaux avantages et les
limites de cette méthode d'analyse.

\begin{table}[H]
\centering
\caption[Avantages et limites des ratios]{Avantages et limites de l'analyse par ratios.}
\label{tab:limitesratios}
\begin{tabularx}{\linewidth}{|X|X|}
\hline
\textbf{Avantages} & \textbf{Limites} \\ \hline
Comparaison entre sociétés de tailles différentes. & Dépendance à la qualité et à la cohérence des états financiers. \\ \hline
Suivi simple de l'évolution dans le temps. & Sensibilité aux méthodes comptables et aux éléments exceptionnels. \\ \hline
Synthèse rapide de la liquidité, de la rentabilité et du risque. & Nécessité d'une interprétation sectorielle. \\ \hline
Base robuste pour les tableaux de bord. & Vision principalement historique, à compléter par la prévision. \\ \hline
\end{tabularx}
\end{table}

%% ───────────────────────────────────────────────────────────────────────────
\section{Référentiel MSI20000 et score de performance}
\label{sec:msi}
\par Le projet s'appuie sur un référentiel industriel normalisé de type
MSI20000. Dans le cadre de ce travail, ce référentiel définit un ensemble de
ratios, des fourchettes de référence et une grille de notation. L'objectif est
de passer d'un ensemble d'indicateurs dispersés à une lecture structurée et
comparable de la performance financière.

\subsection{Logique de notation}
\par Chaque ratio clé reçoit une note comprise entre 0 et 5 selon sa position
par rapport aux seuils retenus. Une note élevée indique que le ratio se situe
dans une zone considérée comme favorable. Les notes sont ensuite agrégées par
catégorie, puis pondérées pour obtenir le \ac{SGPI}. La formule générale peut
s'écrire :

\begin{equation}
SGPI = \sum_{c=1}^{n} \left(\frac{\overline{N_c}}{5} \times P_c\right) \times 100
\end{equation}

\par où $\overline{N_c}$ représente la note moyenne de la catégorie $c$ et
$P_c$ son poids dans le score global. Lorsque certaines catégories ne sont pas
disponibles, les poids utilisés peuvent être renormalisés afin de ne pas
pénaliser mécaniquement une société pour des données absentes. La
figure~\ref{fig:sgpiflow} synthétise cette logique de construction du score.

\begin{figure}[H]
\centering
\begin{tikzpicture}[
  node distance=0.75cm and 0.8cm,
  box/.style={rectangle, rounded corners, draw=darkblue, thick, fill=lightblue!35,
    align=center, minimum width=3.2cm, minimum height=1cm, font=\footnotesize},
  score/.style={rectangle, rounded corners, draw=espritred, very thick,
    fill=espritred!12, align=center, minimum width=3.2cm, minimum height=1cm,
    font=\footnotesize\bfseries},
  arr/.style={-{Latex[length=2.2mm]}, thick, darkblue}
]
\node[box] (r) {Ratios financiers\\46 indicateurs};
\node[box, right=of r] (n) {Notes par ratio\\0 à 5};
\node[box, right=of n] (c) {Moyennes par\\catégorie};
\node[score, right=of c] (s) {SGPI\\/100};
\draw[arr] (r) -- (n);
\draw[arr] (n) -- (c);
\draw[arr] (c) -- (s);
\end{tikzpicture}
\caption[Construction du SGPI]{Construction du Score Global de Performance Industrielle.}
\label{fig:sgpiflow}
\end{figure}

\subsection{Score de Conan \& Holder}
\par Parmi les indicateurs de risque, le \textbf{score de Conan \&
Holder}~\cite{conan1979} occupe une place particulière. Il s'agit d'une fonction
discriminante construite à partir de plusieurs ratios financiers afin d'apprécier
le risque de défaillance. Dans une logique de diagnostic, il complète les ratios
classiques en proposant une lecture synthétique du risque.

\par Le score ne doit pas être interprété isolément. Il doit être rapproché des
autres dimensions : structure financière, rentabilité, liquidité et capacité à
générer des flux. Dans le cadre de ce projet, il est intégré dans la catégorie
de gestion des risques et contribue au \ac{SGPI}.

%% ───────────────────────────────────────────────────────────────────────────
\section{Benchmark financier}
\label{sec:benchmark_theorie}
\par Le benchmark consiste à comparer une entreprise avec des références
pertinentes. En analyse financière, il peut s'agir de concurrents directs,
d'entreprises du même secteur, de leaders internationaux ou de normes
industrielles. L'objectif n'est pas de copier les résultats d'un concurrent,
mais de situer la performance relative d'une entreprise.

\subsection{Conditions d'un benchmark pertinent}
\par Un benchmark financier doit respecter plusieurs conditions. Les indicateurs
comparés doivent être calculés de manière homogène, les périodes doivent être
alignées, les devises et les échelles doivent être identifiées, et les
différences de modèle économique doivent être prises en compte. Les ratios sont
particulièrement utiles car ils réduisent l'effet de taille et de devise, mais
ils ne suppriment pas toutes les différences sectorielles. Le
tableau~\ref{tab:benchmark} présente les conditions nécessaires à un benchmark
fiable.

\begin{table}[H]
\centering
\caption[Conditions d'un benchmark financier]{Conditions d'un benchmark financier fiable.}
\label{tab:benchmark}
\begin{tabularx}{\linewidth}{|p{3.4cm}|X|}
\hline
\textbf{Condition} & \textbf{Justification} \\ \hline
Homogénéité des formules & Deux sociétés ne sont comparables que si les ratios sont calculés de la même manière. \\ \hline
Alignement temporel & Une comparaison doit porter sur la même période ou sur des périodes clairement identifiées. \\ \hline
Prise en compte de la devise & Les montants absolus nécessitent une conversion ; les ratios sont plus directement comparables. \\ \hline
Compréhension du secteur & Une holding, une entreprise industrielle et un groupe mondial n'ont pas exactement le même modèle économique. \\ \hline
Traçabilité des sources & Les écarts doivent pouvoir être expliqués à partir des données d'origine. \\ \hline
\end{tabularx}
\end{table}

\subsection{Choix des sociétés de comparaison}
\par Dans ce projet, la \sfbt est comparée à DELICE, AB~InBev et Coca-Cola. Ce
choix permet de combiner une référence locale et deux références internationales
du secteur des boissons. La comparaison doit toutefois rester prudente : les
tailles, devises, marchés et modèles économiques diffèrent fortement. C'est
précisément pour cette raison que l'analyse privilégie les ratios plutôt que les
montants bruts.

%% ───────────────────────────────────────────────────────────────────────────
\section{Apprentissage automatique et prévision}
\label{sec:ml}
\par L'\textbf{apprentissage automatique} (\ac{ML})~\cite{james2021} regroupe
les méthodes permettant à un modèle d'apprendre une relation à partir de données
historiques. En finance, ces méthodes peuvent servir à prédire des ventes, un
résultat, une trésorerie, un risque ou un score. Dans ce projet, l'objectif est
de prévoir quelques grandeurs financières clés de la \sfbt à partir de leur
historique.

\subsection{Prévision de séries temporelles}
\par Une série temporelle est une suite d'observations ordonnées dans le temps.
Sa prévision consiste à utiliser les valeurs passées pour estimer les valeurs
futures. Les séries financières présentent souvent une tendance, une saisonnalité
et parfois des ruptures. Le modèle doit donc tenir compte de l'ordre temporel :
il serait incorrect de mélanger aléatoirement les observations d'entraînement et
de test comme dans un problème classique de classification.

\par La figure~\ref{fig:forecastflow} montre le principe d'une prévision
supervisée à partir d'une série temporelle : les valeurs passées sont transformées
en variables explicatives, puis le modèle est évalué sur des périodes
postérieures.

\begin{figure}[H]
\centering
\begin{tikzpicture}[
  node distance=0.7cm and 0.9cm,
  box/.style={rectangle, rounded corners, draw=midblue, thick, fill=lightblue!40,
    align=center, minimum width=3cm, minimum height=1cm, font=\footnotesize},
  arr/.style={-{Latex[length=2.2mm]}, thick, midblue}
]
\node[box] (serie) {Historique\\financier};
\node[box, right=of serie] (feat) {Variables\\retards, tendance, saison};
\node[box, right=of feat] (model) {Modèle\\ML};
\node[box, right=of model] (future) {Prévision\\future};
\draw[arr] (serie) -- (feat);
\draw[arr] (feat) -- (model);
\draw[arr] (model) -- (future);
\end{tikzpicture}
\caption[Chaîne de prévision supervisée]{Principe d'une prévision supervisée sur série temporelle.}
\label{fig:forecastflow}
\end{figure}

\subsection{Régression linéaire}
\par La régression linéaire est un modèle statistique qui cherche à expliquer
une variable cible par une combinaison linéaire de variables explicatives. Sa
forme générale est :

\begin{equation}
\hat{y} = \beta_0 + \beta_1 x_1 + \beta_2 x_2 + \cdots + \beta_p x_p
\end{equation}

\par Sa principale force est son interprétabilité. Elle est adaptée lorsque la
relation entre les variables suit une tendance régulière. Dans le cas de séries
financières relativement stables, elle constitue une base solide et défendable.
Sa limite est sa difficulté à capturer des relations fortement non linéaires ou
des interactions complexes si celles-ci ne sont pas explicitement modélisées.

\subsection{Forêt aléatoire}
\par La forêt aléatoire, ou \textit{Random Forest}, est un modèle ensembliste
introduit par Breiman~\cite{breiman2001}. Elle combine plusieurs arbres de
décision entraînés sur des sous-échantillons des données et des sous-ensembles de
variables. La prédiction finale correspond à la moyenne des prédictions des
arbres en régression.

\par Ce modèle est robuste et capte des relations non linéaires. Il nécessite
peu d'hypothèses statistiques et résiste relativement bien au bruit. En revanche,
il est moins interprétable qu'une régression linéaire et peut rencontrer une
limite importante en prévision de séries croissantes : les arbres extrapolent
mal au-delà des niveaux observés. Pour cette raison, le projet prévoit de
modéliser un taux de croissance plutôt que le niveau brut. Le
tableau~\ref{tab:modelesml} compare les deux modèles retenus.

\begin{table}[H]
\centering
\caption[Comparaison des modèles de prévision]{Comparaison des deux modèles de prévision retenus.}
\label{tab:modelesml}
\begin{tabularx}{\linewidth}{|p{3.2cm}|X|X|}
\hline
\textbf{Modèle} & \textbf{Forces} & \textbf{Limites} \\ \hline
Régression linéaire & Simple, interprétable, efficace sur des tendances régulières, facile à défendre dans un rapport. & Capte mal les relations très non linéaires si les variables ne sont pas adaptées. \\ \hline
Random Forest & Robuste, non linéaire, capable de gérer des interactions entre variables. & Moins interprétable, extrapolation limitée sur des séries fortement croissantes. \\ \hline
\end{tabularx}
\end{table}

\subsection{Évaluation des modèles}
\par La qualité d'un modèle de prévision se mesure par plusieurs métriques. Le
coefficient de détermination $R^2$ mesure la part de variance expliquée par le
modèle. La \ac{MAE} mesure l'erreur moyenne absolue. La \ac{RMSE} pénalise plus
fortement les grandes erreurs. La \ac{MAPE} exprime l'erreur en pourcentage, ce
qui facilite l'interprétation métier. Le tableau~\ref{tab:metriques} présente
les métriques utilisées.

\begin{table}[H]
\centering
\caption[Métriques d'évaluation des modèles]{Métriques d'évaluation utilisées pour la prévision.}
\label{tab:metriques}
\begin{tabularx}{\linewidth}{|p{2.5cm}|p{4.2cm}|X|}
\hline
\textbf{Métrique} & \textbf{Formule simplifiée} & \textbf{Interprétation} \\ \hline
MAE & $\frac{1}{n}\sum |y_i-\hat{y_i}|$ & Erreur moyenne en valeur absolue, dans l'unité de la variable. \\ \hline
RMSE & $\sqrt{\frac{1}{n}\sum (y_i-\hat{y_i})^2}$ & Erreur quadratique moyenne, plus sensible aux grandes erreurs. \\ \hline
MAPE & $\frac{100}{n}\sum |\frac{y_i-\hat{y_i}}{y_i}|$ & Erreur moyenne en pourcentage, facilement lisible par un décideur. \\ \hline
$R^2$ & $1-\frac{\sum(y_i-\hat{y_i})^2}{\sum(y_i-\bar{y})^2}$ & Part de variance expliquée ; plus la valeur est proche de 1, meilleur est le modèle. \\ \hline
\end{tabularx}
\end{table}

\par Dans le contexte d'une série temporelle, l'évaluation doit respecter
l'ordre chronologique. Le modèle est entraîné sur le passé puis testé sur une
période postérieure. Cette démarche, appelée \textit{backtest}, reproduit mieux
les conditions réelles de prévision.

%% ───────────────────────────────────────────────────────────────────────────
\section{Synthèse des choix méthodologiques}
\label{sec:synthese_chap2}
\par Les concepts présentés dans ce chapitre guident directement la conception
du projet. Le tableau~\ref{tab:synthesechap2} relie chaque notion théorique à
son application concrète dans la plateforme développée.

\begin{table}[H]
\centering
\caption[Synthèse des choix méthodologiques]{Lien entre état de l'art et choix du projet.}
\label{tab:synthesechap2}
\begin{tabularx}{\linewidth}{|p{3.4cm}|X|X|}
\hline
\textbf{Notion théorique} & \textbf{Principe retenu} & \textbf{Application dans le projet} \\ \hline
Business Intelligence & Transformer la donnée brute en information décisionnelle. & Pipeline automatisé, données préparées, rapports et dashboards Power BI. \\ \hline
ETL & Extraire, transformer et charger des données fiables. & Scripts Python de lecture Excel, nettoyage, normalisation et export. \\ \hline
Qualité des données & Contrôler la cohérence et documenter les anomalies. & Rapport qualité, collisions, contrôle Actif = Passif, tests. \\ \hline
Schéma en étoile & Structurer les données autour de faits et dimensions. & Data Warehouse SQLite avec dimensions société, temps, indicateur et ratio. \\ \hline
Ratios financiers & Comparer la performance indépendamment de la taille. & Calcul de 46 ratios pour SFBT et sociétés de benchmark. \\ \hline
Scoring & Synthétiser plusieurs indicateurs en score. & Score Global de Performance Industrielle sur 100. \\ \hline
Prévision ML & Exploiter l'historique pour anticiper les valeurs futures. & Régression linéaire et Random Forest évalués par backtest. \\ \hline
\end{tabularx}
\end{table}

%% ── Conclusion ─────────────────────────────────────────────────────────────
\section*{Conclusion}
\addcontentsline{toc}{section}{Conclusion}
\par Ce chapitre a posé le socle théorique du projet. La \ac{BI} apporte le
cadre général de transformation des données en information décisionnelle. Le
processus \ac{ETL} et les principes de qualité assurent la fiabilité des données
financières. La modélisation dimensionnelle justifie l'organisation en schéma en
étoile. L'analyse par ratios et le référentiel MSI20000 structurent le
diagnostic financier, tandis que le benchmark situe la \sfbt par rapport à des
références sectorielles. Enfin, l'apprentissage automatique complète l'analyse
par une dimension prédictive.

\par Ces fondements expliquent les choix de conception détaillés dans le
chapitre suivant : architecture de la solution, schéma de données, règles
de traitement, moteur de ratios, score global et module de prévision.
